%% hopfion-framework/foundations/alpha.tex
%% Sources: main_paper3.tex, main_paper4.tex
%% The fine structure constant: two-, three-, and four-term formulas
%% Auto-generated from project files. Edit source papers to update.
%%
%% Status tags: \status{published}, \status{open}, \status{conjectured}
%% \source{PaperN, label} — where the proof lives
%% \residual{X\%} or \residual{exact} — numerical accuracy


%════════════════════════════════════════════════════════════════════
% WZW and Two-Term Alpha — Paper III
%════════════════════════════════════════════════════════════════════
\begin{equation}
  \aInv_{\mathrm{geo}}
  \;=\; \frac{\vp^6}{\sin^4(\pi/5)}
  \;=\; 150.332\ldots
  \label{P3:eq:alpha_geo}
\end{equation}

\begin{equation}
  \aInv(\mu_{\mathrm{IR}})
  \;=\; \aInv(\mu_{\mathrm{UV}})
        - \frac{k}{2\pi}\,\log\frac{\mu_{\mathrm{UV}}}{\mu_{\mathrm{IR}}},
  \label{P3:eq:wzw_running}
\end{equation}

\begin{equation}
  \boxed{\aInv \;=\; \frac{360}{\vp^2} - \frac{k}{2\pi}
         \;=\; 137.030\ldots\quad(\text{gap: }+0.004\%)}
  \label{P3:eq:alpha_formula}
\end{equation}

\begin{equation}
  \delta_\alpha = \frac{\aInv_{\mathrm{formula}} - \aInv_{\mathrm{measured}}}
                       {\aInv_{\mathrm{measured}}}
                = +0.004\%,
  \label{P3:eq:alpha_gap}
\end{equation}

\begin{equation}
  \mu\,\frac{\partial}{\partial\mu}\,\frac{1}{g^2}
  \;=\; -\frac{k}{4\pi^2}\,(1 + O(g^2))
  \quad\Longrightarrow\quad
  \frac{d\,\aInv}{d\log\mu} = -\frac{k}{2\pi}.
  \label{P3:eq:beta_wzw}
\end{equation}

\begin{equation}
  \kappa_{\mathrm{tube}}
  = (f')^2 + \frac{\sin^2\!f}{\rho^2}
  = 2(f')^2
  = g_{\mathrm{WZW}},
  \label{P3:eq:kappa_wzw}
\end{equation}


\begin{result}[Two-step cascade]
\label{P3:prop:cascade}
The inverse fine structure constant exhibits a two-step geometric structure:
\begin{equation}
  \underbrace{\frac{\vp^6}{\sin^4(\pi/5)}}_{150.332}
  \;\xrightarrow{\Delta_1 = -12.824}\;
  \underbrace{\frac{360}{\vp^2}}_{137.508}
  \;\xrightarrow{\Delta_2 = -0.472}\;
  \underbrace{\aInv_{\mathrm{phys}}}_{137.036}
\end{equation}
where $\Delta_1$ is WZW anti-screening from the UV condensate scale to
the golden-angle attractor, and $\Delta_2 = k/(2\pi) = 3/(2\pi)$ is
the one-loop WZW correction.  Derivation of $\Delta_1$ from first
principles is Open Problem~4 (resolved in Paper~IV).
\end{result}
\status{proved (cascade structure; $\Delta_1$ derivation in Paper~IV)}
\source{P3:prop:cascade}
\residual{leading order, 0.004\%}

\begin{result}[Fine structure constant prediction]
\label{P3:res:alpha}
The icosahedral WZW framework predicts
\begin{equation}
  \boxed{\aInv = \frac{360}{\vp^2} - \frac{k}{2\pi} = 137.030\ldots
         \quad(\text{gap: }+0.004\%)}
\end{equation}
where $k=3$ is forced by the $2I$ structure and $360/\vp^2$ is the
golden angle.  The residual $0.006$ is a genuine formula gap (not a
lattice artefact), resolved to $5\times10^{-7}\%$ in Paper~IV.
\end{result}
\status{proved (leading order); improved to $5\times10^{-7}\%$ in Paper~IV}
\source{P3:res:alpha}
\residual{$+0.004\%$; four-term formula achieves $5\times10^{-7}\%$ (Paper~IV)}

%════════════════════════════════════════════════════════════════════
% Three- and Four-Term Formulas — Paper IV
%════════════════════════════════════════════════════════════════════
\begin{remark}[The precise bare-mass condition in 3D]
\label{P4:rem:bare_mass}
Theorem~P4:thm:mass\_renorm
%\ref{P4:thm:mass_renorm}
(filed in \texttt{electroweak/lepton\_masses.tex})
is an exact statement about $\mWZW/m_0$
at the physical torus radius $R_0$.  For the QED running to yield the
observed $\aInv$, we require $\mWZW > m_e$, hence
$m_0 > m_e/0.9042 = 1.106\,m_e$ (at $R_0=3$; compare $1.076\,m_e$ in the
thin-torus limit).
The three-term formula holds if and only if
$\mWZW = m_e\,e^{\pi/(k\vp^6)}$, i.e.\
\begin{equation}
  m_0 = m_e\,\sqrt{\frac{95}{48\vp}}\,e^{\pi/(k\vp^6)} \approx 1.172\,m_e
  \quad (R_0=3),
  \label{P4:eq:bare_mass_3D}
\end{equation}
replacing the thin-torus value $m_0\approx 1.141\,m_e$.
In general:
\[
  m_0 = m_e\,\sqrt{\frac{2R_0^2+1}{2R_0^2}\cdot\frac{15}{8\vp}}\,e^{\pi/(k\vp^6)}.
\]
The exact 3D formula requires a $17\%$ departure from $m_0=m_e$.
The ratio $\mWZW/m_e = e^{\pi/(k\vp^6)}$ is derived from the condensate
geometry to $0.22\%$ accuracy via Theorem~\ref{P4:thm:t_matrix_spoke}; the
analytical derivation of the exact value remains as a higher-order problem.
\end{remark}
\status{published}
\source{P4:rem:bare_mass}

\begin{remark}[Path to single-input: $m_e$ derivable from $\TCMB$ up to $0.22\%$]
\label{P4:rem:single_input}
Theorem~\ref{P4:thm:verlinde_cancel} combined with the spoke formula
establishes that the electron mass is derivable from $\TCMB$ and
$\vp$ alone:
\begin{equation}
  m_e
  = \TCMB\!\left(\frac{\pi^2}{150}\right)^{\!\!1/4}
    \cdot\vp^{20}\cdot
    \exp\!\left(\frac{1600\pi - 3}{400}\right)
    \cdot\bigl(1 + O(R_0^{-2})\bigr).
  \label{P4:eq:me_single_input}
\end{equation}
Every factor is a topological or WZW-algebraic quantity; no Standard
Model parameter enters.
The $O(R_0^{-2})\approx0.22\%$ residual is the $2D$ geometric correction from
the thick-torus toroidal geometry at the physical condensate radius $R_0=3$.
In the thin-torus limit ($R_0\to\infty$) this correction vanishes exactly,
and~\eqref{P4:eq:me_single_input} becomes exact.
Closing the $0.22\%$ analytically, equivalent to proving the normalisation
conjecture $\vp^9\sqrt{\Ja\Jfour}=Q$ of Paper~II at the exact thick-torus
profile, is resolved in Paper~XII (see \texttt{cor:vew}).
The ratio $m_e/\TCMB \approx 2.18\times10^9$ (natural units, $k_B=1$) is a
nearly-determined dimensionless invariant, accurate to $0.22\%$ without any
free parameter.
\end{remark}
\status{published}
\source{P4:rem:single_input}
\residual{$0.22\%$ at physical $R_0=3$; exact in the thin-torus limit}

\begin{theorem}[Three-term fine structure constant (unconditional)]
\label{P4:thm:three_term_alpha}
The fine structure constant at the electron mass scale is:
\begin{equation}
  \boxed{
  \aInv
  = \underbrace{\frac{360}{\vp^2}}_{\text{golden angle}}
  - \underbrace{\frac{k}{2\pi}}_{\text{WZW anti-screening}}
  + \underbrace{\frac{1}{3k\vp^6}}_{\text{QED screening}}
  = 137.036\,49\ldots
  }
  \label{P4:eq:three_term}
\end{equation}
with $k=3$, $\vp^6 = \lambda \approx 17.944$ (the Bogomolny parameter,
Corollary~5.10 of~\cite{Paper1}).  Measured: $\aInv = 137.035\,999\,177$
(PDG~\cite{PDG2024}).  Residual vs measured: $|137.036\,49 - 137.035\,999| = 0.000\,49$,
i.e.\ $0.0004\%$.  This formula is proved unconditionally in
Theorem~\ref{P4:thm:alpha_unconditional} (Section~\ref{sec:unconditional_alpha})
using the modular $T$-matrix correction of Theorem~\ref{P4:thm:t_matrix_spoke}.
\end{theorem}
\status{published}
\source{P4:thm:three_term_alpha}
\residual{$+0.0004\%$}

\begin{remark}[Sign pattern, convergence, and the anti-screening mechanism]
\label{P4:rem:sign_pattern}
The three corrections in Theorem~\ref{P4:thm:three_term_alpha} have
alternating signs: $360/\vp^2=+137.508$ (geometric UV attractor),
$-k/(2\pi)=-0.477$ (WZW anti-screening, non-Abelian), and
$+1/(9\vp^6)=+0.006$ (one-loop QED screening, Abelian $e^+e^-$).
The dominant mechanism is \emph{WZW anti-screening}: the $\mathrm{SU}(2)_3$
condensate is non-Abelian and asymptotically free, in exact analogy with
QCD; virtual flux-knot pairs in the condensate vacuum enhance the coupling
at lower energies, reducing $\aInv$ from the UV geometric value $150.3$
toward the golden-angle attractor $360/\vp^2=137.508$.
The subleading mechanism is ordinary QED screening: virtual $e^+e^-$ pairs
screen the bare charge, giving $\delta_{\rm QED}=+1/(9\vp^6)$ when running
from $\mWZW$ down to $m_e$ --- opposite in sign to the WZW anti-screening,
as expected for a non-Abelian vs.\ Abelian competition.
The residual $0.000\,492$ is closed by the Chern--Simons correction of
Theorem~\ref{P4:thm:four_term_alpha} (a fourth term $-1/(36\pi\vp^6)=-T_3/(4\pi)$,
reducing the residual to $5\times10^{-7}\%$); it is not at the two-loop
QED level $\alpha^2/(3\pi^2)$ ($\sim10^{-5}$, fifty times smaller), and the
WZW T-matrix expansion is complete at three terms since
$\cos(2\pi Q T_{j=1})=\cos(2\pi Q T_{j=3/2})=0$ exactly.
\end{remark}
\status{published}
\source{P4:rem:sign_pattern}

\begin{theorem}[Four-term fine structure constant, $5\times10^{-7}\%$ accuracy]
\label{P4:thm:four_term_alpha}
The fine structure constant at the electron mass scale satisfies
\begin{equation}
  \boxed{
  \aInv
  = \underbrace{\frac{360}{\vp^2}}_{\text{golden angle}}
  - \underbrace{\frac{k}{2\pi}}_{\text{WZW anti-screen.}}
  + \underbrace{\frac{1}{9\vp^6}}_{\text{QED one-loop}}
  - \underbrace{\frac{1}{36\pi\vp^6}}_{\text{CS correction}}
  = 137.035\,998\,486\ldots
  }
  \label{P4:eq:four_term}
\end{equation}
with $k=3$.  The measured value is $\aInv = 137.035\,999\,177$
(PDG~\cite{PDG2024}).  Residual: $|137.035\,998\,486 - 137.035\,999\,177|
= 6.91\times10^{-7}$ in $\aInv$ units
($5\times10^{-7}\%$ fractional accuracy).
The four-term formula represents a factor-of-700 improvement over the
three-term result of Theorem~\ref{P4:thm:three_term_alpha}.
\end{theorem}
\status{published}
\source{P4:thm:four_term_alpha}
\residual{$5\times10^{-7}\%$}

\begin{remark}[Physical content of the Chern--Simons correction]
\label{P4:rem:cs_correction}
The fourth term has a topological interpretation: the WZW condensate
carries a $\mathrm{U}(1)$ magnetic flux $\Phi_{\mathrm{CS}}=k\Phi_0/(4\pi)$
through the condensate torus. When a virtual $e^+e^-$ pair (the QED loop
of Term~3) encloses this flux, its contribution to the vacuum polarisation
is suppressed by the Aharonov--Bohm factor $\cos(k/(4\pi))\approx1-k/(4\pi)$
at small CS coupling; the relevant suppression per unit WZW level is
$1/(4\pi)$, giving $1/(36\pi\vp^6)=1/(9\vp^6)\times1/(4\pi)$. This is
consistent with the topological origin of Term~3 itself
($1/9=Q/(k\,h(E_8))$, Theorem~\ref{P4:thm:three_term_alpha}): both arise from
the WZW/McKay algebraic structure of the condensate, not perturbative QED.
\end{remark}
\status{published}
\source{P4:rem:cs_correction}

\begin{remark}[The four-term series and convergence]
\label{P4:rem:four_term_series}
The four terms of Theorem~\ref{P4:thm:four_term_alpha} form a convergent
series with alternating signs $+,-,+,-$ and decreasing magnitudes:
$360/\vp^2=+137.507\,764$, $-k/(2\pi)=-0.477\,465$,
$+1/(9\vp^6)=+0.006\,192$, $-1/(36\pi\vp^6)=-0.000\,493$, summing to
$137.035\,998$ vs.\ measured $137.035\,999$ ($6.9\times10^{-7}$ absolute,
$5\times10^{-7}\%$ fractional). The ratio of successive terms is
$T_4/T_3\approx1/(4\pi)\approx0.0796$, $T_3/T_2\approx0.013$,
$T_2/T_1\approx0.0035$: the series converges geometrically and is
effectively complete at four terms, the next correction being
$\sim6\times10^{-6}$.
\end{remark}
\status{published}
\source{P4:rem:four_term_series}

\begin{theorem}[Three generations from $\mathrm{SU}(2)_3$]
\label{P4:thm:three_gen}
The $\mathrm{SU}(2)_k$ WZW model at level $k=3$ has exactly three
particle primary fields (spokes), corresponding to the three
Standard-Model fermion generations.
\end{theorem}
\status{published}
\source{P4:thm:three_gen}
\residual{exact}

\begin{theorem}[Spoke ground state]
\label{P4:thm:spoke_ground}
Two first-generation ($j=\tfrac{1}{2}$) spokes interact via the
$\mathrm{SU}(2)_3$ fusion rule~\cite{Verlinde1988}
\begin{equation}
  j_{1/2} \times j_{1/2} = j_0 \oplus j_1,
  \label{P4:eq:fusion}
\end{equation}
producing a singlet ($j=0$) and a triplet ($j=1$) channel.
The $j=0$ channel is the $Q=2$ Hopfion condensate proved in~\cite{Paper3}.
It is the unique ground state: the singlet has conformal dimension
$h_0 = 0$, strictly lower than the triplet $h_1 = j(j+1)/(k+2) = 2/5$.
\end{theorem}
\status{published}
\source{P4:thm:spoke_ground}
\residual{exact}

\begin{theorem}[Spoke mass spectrum]
\label{P4:thm:spoke_mass_spectrum}
Each of the three particle primaries $j\in\{\tfrac{1}{2},1,\tfrac{3}{2}\}$
of $\mathrm{SU}(2)_3$ defines a stable boundary state $|B_j\rangle$ whose
mass satisfies
\begin{equation}
  \frac{m_j}{\Lcond}
  = \vp^{2Q}\cdot\exp\!\left(4\pi - \frac{T_j}{Q}\right),
  \qquad
  T_j = h_j - \frac{c}{24} = \frac{j(j+1)}{k+2} - \frac{c}{24}.
  \label{P4:eq:spoke_mass_spectrum}
\end{equation}
The masses are strictly ordered: $m_{1/2} > m_1 > m_{3/2}$.
\end{theorem}
\status{published}
\source{P4:thm:spoke_mass_spectrum}

\begin{theorem}[Electron is the lightest spoke: derivation without assumption]
\label{P4:thm:me_lightest}
The lightest particle in the $\mathrm{SU}(2)_3$ spoke spectrum is the
$j=\tfrac{1}{2}$ primary.  Its mass equals $m_e$; no external assumption
about $m_e$ enters.
\end{theorem}
\status{published}
\source{P4:thm:me_lightest}
\residual{exact}

\begin{theorem}[Casimir--T-matrix self-duality uniquely selects $j=1/2$]
\label{P4:thm:casimir_selects}
The k=3 self-duality $T_{1/2} = c/24 = 3/40$ (proved in the companion
quantum-correction note) has the following consequence: the lightest
particle T-matrix phase equals the WZW Casimir energy coefficient.
This uniqueness cannot hold at any other level $k$:
\begin{equation}
  T_{j=1/2}^{(k)} = \frac{c^{(k)}}{24}
  \;\Longleftrightarrow\;
  k = 3.
  \label{P4:eq:k3_unique}
\end{equation}
\end{theorem}
\status{published}
\source{P4:thm:casimir_selects}
\residual{exact}

\begin{theorem}[Lower segment of the first-generation spoke]
\label{P4:thm:second_spoke}
The ratio of the condensate scale to the electron mass satisfies
\begin{equation}
  \frac{m_e}{\Lcond} = \vp^{2Q}\cdot e^{4\pi}\cdot(1 + O(R_0^{-2}))
  = \vp^{20}\cdot e^{4\pi}\cdot(1 + O(R_0^{-2})),
  \label{P4:eq:spoke_leading}
\end{equation}
where the $O(R_0^{-2})$ correction amounts to $0.97\%$ at $R_0=3$.
The $\vp^{20}$ factor is exact (Verlinde + $Q=10$ topology).
The $e^{4\pi}$ factor is the leading-order Euclidean CS contribution
$e^{-\Delta S_\mathrm{CS}}$ from the Abelian Hopf Chern-Simons action difference
between the $Q=2$ and $Q=1$ sectors (equation~\eqref{P7:eq:delta_cs}).
The identification of the Euclidean path-integral ratio with $e^{-\Delta S_\mathrm{CS}}$
is a leading-order heuristic; the precise statement is that the
topological contribution to $\log(m_e/\Lcond)$ from the CS action is $4\pi$.
\end{theorem}
\status{published}
\source{P4:thm:second_spoke}
\residual{$0.97\%$ at $R_0=3$}

\begin{theorem}[Modular $T$-matrix correction to the spoke]
\label{P4:thm:t_matrix_spoke}
The Abelian Chern-Simons spoke exponent $4\pi$ receives a non-Abelian
correction from the modular $T$-matrix phase accumulated by the $j=1/2$
spoke carrier as it traverses the $Q=10$ topological sectors of the
condensate.  The corrected formula is
\begin{equation}
  \boxed{
  \frac{m_e}{\Lcond}
  = \vp^{20}\,\exp\!\Bigl(4\pi - \frac{T_{1/2}}{Q}\Bigr)
    + O(R_0^{-2})
  = \vp^{20}\,\exp\!\Bigl(4\pi - \frac{3}{400}\Bigr)
    + O(R_0^{-2}).
  }
  \label{P4:eq:spoke_corrected}
\end{equation}
Numerically, $\vp^{20}e^{4\pi-3/400} = 4.305\times10^9$ vs.\ measured
$4.296\times10^9$: a residual of $+0.22\%$, a factor-of-$4.4$ improvement.
The $+0.22\%$ residual is identified in Theorem~\ref{P4:thm:msbar_pole}
as the one-loop QED pole-mass to $\overline{\mathrm{MS}}$ scheme conversion
$\aem/\pi = 0.2323\%$, which closes the residual to $0.013\%$.
In the thin-torus limit the genuine gap reduces from $+0.71\%$ to $+0.04\%$,
a factor-of-$18$ improvement.
\end{theorem}
\status{published}
\source{P4:thm:t_matrix_spoke}
\residual{$+0.22\%$}

\begin{remark}[Revised residual decomposition after Verlinde cancellation]
\label{P4:rem:revised_decomp}
Theorem~\ref{P4:thm:verlinde_cancel} establishes that the Verlinde
finite-torus correction is algebraically zero. The full $0.97\%$
leading-order residual of Theorem~\ref{P4:thm:second_spoke} is therefore
entirely genuine, explained by the $T$-matrix correction ($-0.75\%$,
Theorem~\ref{P4:thm:t_matrix_spoke}) plus the $O(R_0^{-2})$ 2D geometry
correction ($+0.22\%$ at physical $R_0=3$; $+0.04\%$ in the thin-torus
limit $R_0\to\infty$). The $+0.22\%$ at $R_0=3$ is a physical
$O(R_0^{-2})$ correction from the full 2D toroidal geometry, not a
numerical grid artefact: Papers~I--III prove $V^*=\vp$ and
$J_4/J_{2a}=2^{4/3}/\vp^5$ exactly at all $R_0$, independently of grid
spacing, with the small $O(h^2)$ solver deviations ($0.013\%$ in $V^*$,
$0.020\%$ in $J_4/J_{2a}$) vanishing as $h\to0$ while the $0.22\%$
physical residual persists.
\end{remark}
\status{published}
\source{P4:rem:revised_decomp}

\begin{theorem}[Unconditional three-term $\alpha^{-1}$ formula]
\label{P4:thm:alpha_unconditional}
The three-term fine structure constant formula
\begin{equation}
  \boxed{
  \aInv = \frac{360}{\vp^2} - \frac{k}{2\pi} + \frac{1}{3k\vp^6}
        = \frac{360}{\vp^2} - \frac{3}{2\pi} + \frac{1}{9\vp^6}
        = 137.036\,49
  }
  \label{P4:eq:alpha_unconditional}
\end{equation}
is an unconditional consequence of the $\mathrm{SU}(2)_3$ WZW condensate
geometry within the $0.22\%$ accuracy of Theorem~\ref{P4:thm:t_matrix_spoke}.
The gap from the measured value $\aInv_{\rm meas}=137.035\,999\,177$
is $+0.000\,49$ ($+0.0004\%$).
\end{theorem}
\status{published}
\source{P4:thm:alpha_unconditional}
\residual{$+0.0004\%$}

\begin{theorem}[QED pole-mass to $\overline{\mathrm{MS}}$ correction]
\label{P4:thm:msbar_pole}
The spoke formula (Theorem~\ref{P4:thm:t_matrix_spoke}) computes
$m_e/\Lcond$ in the $\overline{\mathrm{MS}}$ renormalisation scheme,
native to the WZW bosonization dictionary.
The measured electron mass $m_e = 0.51100\,\mathrm{MeV}$ is the
\emph{pole mass}. The one-loop QED conversion is:
\begin{equation}
  m_e^{\mathrm{pole}} = m_e^{\overline{\mathrm{MS}}}(m_e)\,
  \Bigl(1 + \frac{\aem}{\pi}\Bigr),
  \label{P4:eq:pole_msbar}
\end{equation}
so that the scheme-corrected spoke formula reads:
\begin{equation}
  \boxed{
  \frac{m_e}{\Lcond} = \vp^{20}\,\exp\!\Bigl(
    4\pi - \frac{3}{400} - \frac{\aem}{\pi}\Bigr)
  = 4.296\times10^9,
  }
  \label{P4:eq:spoke_pole_corrected}
\end{equation}
giving a residual of $-0.013\%$ vs.\ the measured $4.296\times10^9$.

The three exponent terms have distinct physical origins:
\begin{itemize}[noitemsep]
  \item $4\pi$: one-loop Chern--Simons spoke amplitude (exact, topological);
  \item $-3/400 = -T_{1/2}/Q$: modular $T$-matrix phase of the
    $j=1/2$ spoke carrier (exact, WZW);
  \item $-\aem/\pi$: one-loop QED $\overline{\mathrm{MS}}$-to-pole-mass
    scheme conversion (standard QED, one-loop exact).
\end{itemize}
The same $\aem$ that appears in the four-term $\alpha^{-1}$ formula
(Theorem~\ref{P4:thm:four_term_alpha}) closes the $m_e$ spoke residual:
$\aem/\pi = 0.2323\%$ matches the old residual $+0.22\%$ to $0.01\%$.
The two results are structurally linked through the Chern--Simons
geometry of the condensate background.
\end{theorem}
\status{published}
\source{P4:thm:msbar_pole}
\residual{$-0.013\%$}

\begin{remark}[As $\alpha$ refines, so does $m_e$]
\label{P4:rem:alpha_me_tracking}
The corrected spoke formula (Theorem~\ref{P4:thm:msbar_pole}) establishes a
precise quantitative link between $\aem$ and $m_e$ with no analogue in
the Standard Model. Since $\Lcond$ is fixed by $\TCMB$ alone, $m_e$ is a
function of $\aem$ and nothing else, Eq.~\eqref{P4:eq:me_alpha_sensitivity}.
The amplification factor $1/(\pi\aem)=\aInv/\pi\approx43.6$ means the
fractional shift in $m_e$ is 43.6 times the fractional shift in $\aem$,
giving a zero-parameter, testable consistency condition among the three
independently measured constants $\aem$, $m_e$, $\TCMB$. At the current
best measurement of $\aem$ from the electron anomalous magnetic moment
($\delta\aem/\aem\approx2\times10^{-10}$), this predicts a corresponding
$m_e$ shift of $\approx8.7\times10^{-9}$, comparable to current
Penning-trap precision on $m_e$.
\end{remark}
\status{published}
\source{P4:rem:alpha_me_tracking}

\begin{remark}[Why WZW bosonization gives the $\overline{\mathrm{MS}}$ mass]
\label{P4:rem:msbar_origin}
Witten non-abelian bosonization identifies the WZW fermion mass with the
$\overline{\mathrm{MS}}$ mass at the bosonization scale, because the
bosonization dictionary is derived from the path integral with dimensional
regularisation and $\overline{\mathrm{MS}}$ subtraction. The pole mass
differs by the on-shell self-energy $\delta m/m=\aem/\pi$ at one loop,
independent of any cutoff. The $+0.22\%$ residual of the three-term spoke
was therefore not a geometric failure but a renormalisation-scheme
effect: the framework predicted the correct $\overline{\mathrm{MS}}$ mass,
and the remaining gap was the scheme conversion.
\end{remark}
\status{published}
\source{P4:rem:msbar_origin}

\begin{proposition}[Algebraic status of $b^*$]
\label{P4:prop:bstar}
The self-consistent coupling $b^*$ is the unique root of the
transcendental equation
\begin{equation}
  \frac{15}{8}\,\frac{\arctan\!\sqrt{8b}}{\sqrt{8b}} = \vp,
  \label{P4:eq:bstar_eq}
\end{equation}
and is not expressible in closed $\vp$-algebraic form.  The approximation
$b^* \approx 3\vp^6/800$ is accurate to $0.21\%$.  To next order:
\begin{equation}
  b^* = \frac{3\vp^6}{800} - \frac{V(3\vp^6/800) - \vp}{V'(3\vp^6/800)}
  + O\!\left(\bigl(V(b_0)-\vp\bigr)^2\right),
  \label{P4:eq:bstar_newton}
\end{equation}
where $V'(b) = \tfrac{15}{2}/[8b(1{+}8b)] - 4\vp/(8b)$ evaluated at the
fixed point, and the corrected value is accurate to $0.09\%$.
No $\vp$-algebraic expression for $b^*$ with error below $0.03\%$
shorter than $\vp^n/m$ for large $n$ is known; the closest is $\vp^2/39$
(error $0.033\%$), which lacks a natural geometric interpretation.
\end{proposition}
\status{published}
\source{P4:prop:bstar}
\residual{$0.21\%$ (leading approximation); $0.09\%$ (next order)}

\begin{proposition}[Decomposition of the spoke exponent]
\label{P4:prop:exponent_decomp}
The spoke exponent $4\pi$ decomposes as
\begin{align}
  4\pi
  &= 2\pi\,t_{\rm open}\,T_1
   = 2\pi\,t_{\rm open}\,h_1
     - 2\pi\,t_{\rm open}\,\frac{c}{24} \notag\\
  &= \underbrace{\frac{64\pi}{13}}_{\text{conformal weight}}
     - \underbrace{\frac{12\pi}{13}}_{\text{Casimir}},
  \label{P4:eq:exponent_decomp}
\end{align}
with $t_{\rm open}=80/13$ (Definition~\ref{P4:def:t_open} below),
$h_1=2/5$, $c/24=3/40$.
\end{proposition}
\status{published}
\source{P4:prop:exponent_decomp}
\residual{exact}

\begin{definition}[Open-string modular parameter]
\label{P4:def:t_open}
$t_{\rm open}$ is fixed by equating the WZW leading amplitude to the
Abelian CS action:
\begin{equation}
  e^{-2\pi t_{\rm open}\,T_1} \stackrel{!}{=} e^{-4\pi}
  \;\Longrightarrow\;
  t_{\rm open} = \frac{2}{T_1} = \frac{80}{13}.
  \label{P4:eq:t_open_val}
\end{equation}
The descendant corrections to $\chi_1$ are $O(e^{-160\pi/13})\approx10^{-17}$,
so this value is essentially exact.
\end{definition}
\status{published}
\source{P4:def:t_open}

\begin{lemma}[Locked Hopf angle]
\label{P4:lem:locked_angle}
For the Battye--Sutcliffe profile $f = 2\arctan(C/\rho)$, the gradient $\nabla f$
makes an exact angle of $\pi/4$ with the radial direction $\hat{e}_\rho$
at every point of the tube cross-section:
\begin{equation}
  \sin^2\!\theta = \frac{1}{2},
  \qquad \theta = \frac{\pi}{4}, \qquad \forall\,\rho > 0.
  \label{P4:eq:locked_angle}
\end{equation}
Equivalently, the winding and tube curvature components satisfy the exact identity
\begin{equation}
  \frac{\sin^2\!f}{\rho^2} = \frac{\kappa_\mathrm{tube}}{2},
  \label{P4:eq:kappa_identity}
\end{equation}
where $\theta$ is defined by $\cos\theta = f'/|\nabla f|$.
\end{lemma}
\status{published}
\source{P4:lem:locked_angle}
\residual{exact}

\begin{lemma}[Winding virial factorisation]
\label{P4:lem:virial_factor}
The winding curvature $\kappa_\mathrm{wind} = \sin^2\!f/r^2$ at the torus
surface ($r \approx R_0$) satisfies
\begin{equation}
  \kappa_\mathrm{wind} \approx \frac{\kappa_\mathrm{tube}}{2R_0^2},
  \label{P4:eq:kappa_wind_exact}
\end{equation}
and consequently the full unrenormalized virial at torus radius $R_0$ factors as
\begin{equation}
  V(0;\,R_0)
  = \frac{J_{2\mathrm{iso}}}{J_{2a}}
  = \frac{15}{8}\!\left(1 + \frac{1}{2R_0^2}\right)
  = V(0;\,\infty)\!\left(1 + \frac{1}{2R_0^2}\right),
  \label{P4:eq:virial_factor}
\end{equation}
where $V(0;\,\infty) = 15/8$ is the thin-torus value.
The winding correction is suppressed by $1/(2R_0^2)$ relative to the tube
contribution and vanishes in the thin-torus limit $R_0\to\infty$.
\end{lemma}
\status{published}
\source{P4:lem:virial_factor}
\residual{exact}

\begin{lemma}[Self-duality identity at $k=3$]
\label{P4:lem:self_dual}
In the $\mathrm{SU}(2)_3$ WZW model at level $k=3$:
\begin{equation}
  h_{1/2} = \frac{j(j+1)}{k+2}\bigg|_{j=1/2} = \frac{3}{20}
  = \frac{c}{12} = \frac{9/5}{12},
  \label{P4:eq:h12_c12}
\end{equation}
and therefore the modular $T$-matrix eigenvalue of the $j=1/2$ primary is
\begin{equation}
  T_{1/2}
  = h_{1/2} - \frac{c}{24}
  = \frac{c}{12} - \frac{c}{24}
  = \frac{c}{24}
  = \frac{3}{40}.
  \label{P4:eq:T12_val}
\end{equation}
This identity is equivalent to the $S$-matrix self-duality
$S_{1/2,1/2}=S_{0,0}$ proved in Theorem~10.4 of~\cite{Paper3}.
\end{lemma}
\status{published}
\source{P4:lem:self_dual}
\residual{exact}

\begin{lemma}[Three equivalent forms of the correction]
\label{P4:lem:three_forms}
The spoke correction satisfies the following exact equalities:
\begin{equation}
  \frac{T_{1/2}}{Q}
  = \frac{h_{1/2}-c/24}{Q}
  = \frac{c/24}{Q}
  = \frac{c}{24\,Q}
  = \frac{9/5}{240}
  = \frac{3}{400},
  \label{P4:eq:three_forms}
\end{equation}
all exact at $k=3$, $Q=10$.  The second equality uses $h_{1/2}=c/12$
(Lemma~\ref{P4:lem:self_dual}); the remaining equalities are arithmetic.
\end{lemma}
\status{published}
\source{P4:lem:three_forms}
\residual{exact}

\begin{corollary}[Thin-torus limit]
\label{P4:cor:thin_torus}
As $R_0\to\infty$, equation~\eqref{P4:eq:mass_ratio} reduces to
$\mWZW^2/m_0^2 = 8\vp/15$, recovering the thin-torus result.
At the physical torus radius $R_0=3$, the correction factor is
\[
  \frac{2R_0^2}{2R_0^2+1}\bigg|_{R_0=3}
  = \frac{18}{19},
\]
so $\mWZW^2/m_0^2 = 8\vp\cdot18/(15\cdot19) = 48\vp/95 \approx 0.8175$.
The condensate softens the fermion mass by a factor
$\sqrt{48\vp/95} \approx 0.9042$
($9.6\%$ softening at $R_0=3$, vs $7.1\%$ in the thin-torus limit).
\end{corollary}
\status{published}
\source{P4:cor:thin_torus}
\residual{exact}

\begin{corollary}[Electroweak scale from $\TCMB$, leading order]
\label{P4:cor:vew}
Combined with the proved Yukawa $y_e = e^{-25\pi/6}$ of~\cite{Paper1}:
\begin{equation}
  \vEW = \Lcond\cdot\vp^{20}\cdot e^{49\pi/6}\cdot(1+O(R_0^{-2})),
  \label{P4:eq:vew_leading}
\end{equation}
accurate to $1.45\%$ at physical $R_0=3$, with no Standard Model parameters.
The corrected formula incorporating the $T$-matrix spoke correction is
$\vEW = \Lcond\cdot\vp^{20}\cdot e^{49\pi/6-3/400}+O(R_0^{-2})$,
accurate to $0.69\%$ (Corollary~\ref{P4:cor:vew_corrected}).
\end{corollary}
\status{published}
\source{P4:cor:vew}
\residual{$1.45\%$}

\begin{result}[Hopf-spoke formula for $\vEW$]
\label{P4:conj:spoke}
The electroweak VEV and the CMB condensate scale satisfy
\begin{equation}
  \vEW
  = \Lcond\cdot\vp^{2Q}\cdot e^{49\pi/6 - 3/400}\cdot(1+O(R_0^{-2})),
  \qquad
  \Lcond = \TCMB\!\left(\frac{\pi^2}{150}\right)^{\!1/4},
  \label{P4:eq:spoke_formula}
\end{equation}
proved in Theorem~\ref{P4:thm:second_spoke} (leading order, $1.45\%$) and
Corollary~\ref{P4:cor:vew_corrected} (with $T$-matrix correction, $0.69\%$).
The Fibonacci identity $\vp^{20} = \sqrt{5}\,F(20)$ (since
$\vp^{20}/F(20) = \vp + 1/\vp = \sqrt{5}$) is exact.
\end{result}
\status{published}
\source{P4:conj:spoke}

\begin{corollary}[Corrected electroweak scale formula]
\label{P4:cor:vew_corrected}
The electroweak vacuum expectation value satisfies
\begin{equation}
  v_{\rm EW}
  = \Lcond\cdot\vp^{20}\cdot\exp\!\Bigl(\frac{49\pi}{6} - \frac{3}{400}\Bigr)
  + O(R_0^{-2}),
  \label{P4:eq:vew_corrected}
\end{equation}
accurate to $0.69\%$ at $R_0=3$ (improved from $1.45\%$).
The QED pole-mass correction $-\aem/\pi$ of Theorem~\ref{P4:thm:msbar_pole}
does not appear here: $v_{\rm EW} = m_e/y_e$ and the scheme conversion
in $m_e$ is cancelled by the identical conversion in the Yukawa coupling
$y_e$, since both are $\overline{\mathrm{MS}}$ quantities and $v_{\rm EW}$
is measured as a tree-level relation via the Fermi constant.
The $0.69\%$ residual is a separate $O(R_0^{-2})$ correction from the
running of $y_e$ between $\Lcond$ and $m_e$.
\end{corollary}
\status{published}
\source{P4:cor:vew_corrected}
\residual{$0.69\%$}

\begin{corollary}[Unconditional four-term $\aInv$ formula]
\label{P4:cor:four_term_unconditional}
Since Theorem~\ref{P4:thm:alpha_unconditional} makes the three-term formula
unconditional, the four-term formula of Theorem~\ref{P4:thm:four_term_alpha}
is also unconditional:
\begin{equation}
  \boxed{
  \aInv
  = \frac{360}{\vp^2} - \frac{k}{2\pi}
  + \frac{1}{9\vp^6} - \frac{1}{36\pi\vp^6}
  = 137.035\,998\,486\ldots
  }
  \label{P4:eq:four_term_unconditional}
\end{equation}
with residual $6.9\times10^{-7}$ ($5\times10^{-7}\%$) from the PDG value
$137.035\,999\,177$, a factor-of-$700$ improvement over the three-term result.
\end{corollary}
\status{published}
\source{P4:cor:four_term_unconditional}
\residual{$5\times10^{-7}\%$}

\begin{equation}
  V(\beta)
  = \frac{J_\mathrm{fb}(\beta)}{J_{2a}}
  = \int_0^\infty \frac{\tilde\rho(m^2)}{1+\beta m^2}\,dm^2,
  \label{P4:eq:V_spectral}
\end{equation}

\begin{equation}
  \boxed{
  \aInv = \frac{360}{\vp^2} - \frac{k}{2\pi} + \frac{1}{3k\vp^6}
        = \frac{360}{\vp^2} - \frac{3}{2\pi} + \frac{1}{9\vp^6}
        = 137.036\,49
  }
\end{equation}

\begin{equation}
  m_0 = m_e\,\sqrt{\frac{95}{48\vp}}\,e^{\pi/(k\vp^6)} \approx 1.172\,m_e
  \quad (R_0=3),
\end{equation}

\begin{equation}
  \frac{\mWZW^2}{m_0^2}
  = \frac{V^*}{V(0;\,R_0)}
  = \frac{\vp}{\dfrac{15}{8}+\dfrac{15}{16R_0^2}}
  = \frac{8\vp}{15}\cdot\frac{2R_0^2}{2R_0^2+1},
\end{equation}

\begin{equation}
  \frac{\delta m_e}{m_e} = -\frac{\delta\aem}{\pi\,\aem}
  \;=\; -\frac{\delta\aem}{\aem}\cdot\frac{1}{\pi\aem}
  \;\approx\; -43.6\;\frac{\delta\aem}{\aem}.
  \label{P4:eq:me_alpha_sensitivity}
\end{equation}

\begin{equation}
  m_e^{\mathrm{pole}} = m_e^{\overline{\mathrm{MS}}}(m_e)\,
  \Bigl(1 + \frac{\aem}{\pi}\Bigr),
\end{equation}

\begin{equation}
  R(\beta) = \int_0^\infty \frac{\rho(m^2)}{1 + \beta m^2}\,dm^2,
  \label{P4:eq:spectral}
\end{equation}

\begin{equation}
  \frac{m_j}{\Lcond}
  = \vp^{2Q}\cdot\exp\!\left(4\pi - \frac{T_j}{Q}\right),
  \qquad
  T_j = h_j - \frac{c}{24} = \frac{j(j+1)}{k+2} - \frac{c}{24}.
\end{equation}

\begin{equation}
  \int_0^\infty \frac{\tilde\rho(m^2)}{1+\bst m^2}\,dm^2 = \vp.
  \label{P4:eq:vfp_spectral}
\end{equation}

\begin{equation}
  V(b) \;=\; \frac{15}{8}\cdot\frac{\arctan\sqrt{8b}}{\sqrt{8b}},
  \qquad b = \bst/C^2,
  \label{P4:eq:virial}
\end{equation}

\begin{equation}
  V(0;\,R_0)
  = \frac{J_{2\mathrm{iso}}}{J_{2a}}
  = \frac{15}{8}\!\left(1 + \frac{1}{2R_0^2}\right)
  = V(0;\,\infty)\!\left(1 + \frac{1}{2R_0^2}\right),
\end{equation}

